\documentclass[11pt]{article}

% "review" = anonymous submission version (line numbers, no author block).
% Change to "final" for camera-ready, or "preprint" for a named preprint with
% page numbers -- build the supervisor's reading copy as "preprint".
\usepackage[review]{acl}

\usepackage{times}
\usepackage{latexsym}
\usepackage[T1]{fontenc}
\usepackage[utf8]{inputenc}
\usepackage{microtype}
\usepackage{graphicx}
\usepackage{booktabs}
\usepackage{amsmath}
\usepackage{amssymb}

\title{When the Evaluator Decides the Result:\
       Measuring Faithfulness in Retrieval-Augmented Generation}

\author{Shakhriyorbek Boltabaev \
  Institute of Informatics, University of Szeged \
  Szeged, Hungary \
  \texttt{shakhriyorbekboltabaev@gmail.com}}

\begin{document}
\maketitle

\begin{abstract}
Faithfulness in retrieval-augmented generation is commonly measured by asking a natural language inference model whether the generated answer is entailed by the retrieved context, aggregated by taking the maximum over retrieved chunks. We hold a generation grid fixed --- four embedding models, two question answering datasets, three generators, 24,000 generated answers --- and vary only the faithfulness evaluator. Across six dataset-generator cells, the conclusion about whether embedding choice affects faithfulness changes in two of them depending on which evaluator is used. We trace this to measured properties of the evaluators rather than to the retrieval systems. Replacing a grounded value in an answer with one absent from the context --- the failure a document question answering system most needs to detect --- moves the entailment score by 0.054 on Natural Questions and 0.085 on HotpotQA for the most verbose generator, against 0.262 and 0.271 for a second and 0.451 and 0.391 for a third, and is caught by a fixed threshold 5-14\% of the time against 35-55\% and 57-71\%. The gap is largely accounted for by answer structure: the score falls as the number of assertions in an answer grows, and the three generators average 3.08, 1.77 and 1.25 assertions per answer. It is not accounted for entirely --- at a matched single assertion they still differ, +0.102 against +0.409 and +0.485 --- so answer structure is a partial account rather than a complete explanation. Conditioning on assertion count groups the open-weight generator with the terser proprietary one and leaves the verbose one as the outlier, a distinction a two-generator design cannot draw. Claim-level aggregation and an alternative evaluator each recover part of the sensitivity but neither removes the effect. We also show that including refusals in the evaluation population produces an apparent relationship between retrieval quality and faithfulness that disappears when they are excluded. We report what these measurement choices do to a concrete equivalence claim, and give the conditions under which such a claim can be stated.
\end{abstract}

\section{Introduction}
\label{sec:intro}

A document question answering system is given a set of uploaded files and must answer from them. If an invoice records a project budget of \$1,000, an answer stating \$1,500 is a failure of a specific kind: it is fluent, on topic, cites the right document, and is wrong in the one place that matters. Systems of this kind are typically guarded by a faithfulness score --- a measure of whether the generated answer is supported by the retrieved text --- with answers below a threshold suppressed or flagged.

This paper asks what such a score actually measures. We do not propose a new evaluator. We take a fixed set of 24,000 generated answers and vary the measurement, and report where the measurement rather than the system determines the conclusion.

Our contributions are:

\begin{itemize}
\item A controlled falsification experiment. Replacing one grounded value in an answer with a value absent from the retrieved context, holding everything else fixed, gives a direct measurement of evaluator sensitivity to the error type that matters (Section~\ref{sec:sensitivity}).
\item Evidence that the sensitivity depends strongly on answer structure, not only on the evaluator: entailment scores fall as the number of assertions per answer grows, for all three generators tested (Section~\ref{sec:assertions}).
\item A demonstration that the choice of evaluation population --- specifically, whether refusals are included --- can produce an apparent relationship between retrieval quality and faithfulness where none survives their exclusion (Section~\ref{sec:population}).
\item A statement of what an equivalence claim about faithfulness requires: naming the evaluator, reporting the margin as a curve rather than a point, and anchoring that margin to a measured effect size (Section~\ref{sec:equivalence}).
\end{itemize}

The embedding-model grid is a testbed, not the object of study. It supplies a set of retrieval systems spanning 4.4 and 11.9 NDCG@5 points on the two datasets, with generation and prompting held constant, which is what makes the measurement comparison controlled.


\section{Related Work}
\label{sec:related}


\subsection{Retrieval quality and downstream outcome}
\label{sec:rq-downstream}

\citet{salemi-zamani-2024} report that human document-level relevance labels, and LLM relevance judgments, correlate only weakly with downstream retrieval-augmented generation performance, and propose eRAG, which labels each retrieved document by the downstream result the system's own generator produces from it alone. Their strongest baseline is answer-containment labelling, the scheme our qrels use, which retains a moderate correlation (Kendall's tau 0.31-0.40 on the question-answering datasets); their reported correlations also decline as the number of retrieved documents grows, and are measured at k = 50 against ours at k = 5. Their downstream measure is correctness throughout, obtained from a fine-tuned 60M-parameter encoder-decoder reader. Our starting point is downstream of theirs: we take as given that retrieval metrics do not directly predict generation quality, and ask which downstream property a given measurement instrument can resolve at all --- for faithfulness rather than correctness, and for prompted instruction-following generators, which can abstain.


\subsection{Faithfulness measurement}
\label{sec:faith-measure}

Entailment-based faithfulness scoring is standard: a natural language inference model is asked whether the retrieved context entails the generated answer. RAGAS \citep{es-etal-2024} packages this alongside context and answer relevance metrics. AlignScore \citep{zha-etal-2023} trains a unified alignment function over several factual-consistency tasks rather than relying on a general-purpose inference model. FaithJudge \citep{tamber-etal-2025} uses a language model as the judge. Our work is orthogonal to proposing another such metric: we measure how three of them respond to a controlled, known error, and report the consequences for conclusions drawn with them.


\subsection{Meta-evaluation of faithfulness detectors}
\label{sec:meta-eval}

A separate line of work evaluates the detectors themselves. TRUE \citep{honovich-etal-2022} standardises factual consistency evaluation over eleven annotated datasets and finds inference-based and question-generation approaches to be strong and complementary. AGGREFACT \citep{tang-etal-2023} aggregates nine summarisation factuality datasets, stratifies them by the model that produced the errors, and reports that no single metric is best in all settings or for all error types. RAGTruth \citep{niu-etal-2024} supplies word-level hallucination annotations over roughly 18,000 retrieval-augmented responses. FaithBench \citep{bao-etal-2025} is assembled from summaries on which current detectors disagree, and reports detection accuracies near 50\% on them. MiniCheck \citep{tang-etal-2024} unifies recent grounding datasets as LLM-AGGREFACT and trains compact fact-checkers against that benchmark.

These works compare detectors against gold hallucination labels on a labelled corpus, and the quantity they report is detector accuracy. Our design contains no gold hallucination label. We hold the generation grid, the queries and the generated answers fixed, swap the evaluator, and ask whether the ranking of the systems under study survives. The two questions come apart: a detector that is more accurate on a labelled corpus need not be one that resolves differences between the systems a practitioner is actually choosing between, and it is the second property that decides whether a published comparison describes the systems or the instrument. \citet{xiao-etal-2023} argue on general grounds that generation metrics should be analysed as measurement instruments with stated reliability and validity. The falsification probe in Section~\ref{sec:sensitivity} and the equivalence curve in Section~\ref{sec:equivalence} are two such analyses, carried out on the instruments a RAG faithfulness result is usually reported with.


\subsection{Hallucination and grounding analysis}
\label{sec:hallu}

ReDeEP \citep{sun-etal-2025} localises hallucination to attention heads and feed-forward components. \citet{sinha-2025} reports that embedding-based detection separates synthetic hallucinations well and real ones poorly, the hard cases being those that remain semantically close to a faithful response. \citet{chen-etal-2025} vary the embedding model in a RAG pipeline and combine several of them for mathematics question answering, on the observation that different embedders succeed on different queries. Our analysis is correlational and behavioural rather than mechanistic: we do not probe model internals.


\section{Definitions}
\label{sec:defs}


\subsection{Faithfulness}
\label{sec:def-faith}

For a generated answer a and retrieved context R = \{d1, ..., dk\}, the standard aggregate is:

\begin{equation}
F_{\max}(a, R) \;=\; \max_{d \in R} \mathrm{NLI}(d \models a)
\end{equation}

We also evaluate a claim-level aggregate. Writing C(a) for the assertions in a:

\begin{equation}
F_{\mathrm{claim}}(a, R) \;=\; \min_{c \in C(a)} \; \max_{d \in R} \mathrm{NLI}(d \models c)
\end{equation}

The maximum over chunks is appropriate --- any retrieved document may support a claim. The maximum over claims, implicit in (1), is not: an answer is only as supported as its least supported assertion, so one unsupported claim should not be masked by several supported ones. Faithfulness is measured against the generated answer and the retrieved context only; the gold answer is never used.


\subsection{Correctness and retrieval hit}
\label{sec:def-correct}

Correctness is measured against the gold answer and never against the context, so that a grounded wrong answer and an ungrounded right answer remain distinguishable. A retrieval hit means the gold answer string appears in the retrieved context. This is a strong definition of retrieval success, and the Limitations section notes what it does not capture.


\subsection{On differencing retrieval quality and faithfulness}
\label{sec:def-gap}

An earlier version of this work defined a gap metric as the difference between NDCG@5 and faithfulness, with a normalised variant dividing by retrieval quality. We report the properties of that family here because they are general rather than specific to our implementation. First, the difference has no common unit: a rank-based retrieval metric and an entailment probability are not commensurable. Second, on our data the difference is negative for every model on Natural Questions, since faithfulness (0.85) exceeds retrieval quality (0.79) --- the metric presumes generation loses fidelity relative to retrieval, and it does not. Third, any metric of the form 1 - F/RQ reproduces the retrieval ranking whenever F is near-constant across systems, so apparent agreement of such a metric across generators is uninformative: retrieval quality is generator-invariant by construction. We therefore report absolute faithfulness throughout. Retiring the metric also disposes of a known objection to it at the root rather than patching it: a difference cannot separate a system whose retrieval and faithfulness are both high from one where both are low, and the normalised variant rescales that ambiguity without removing it. Reporting the two quantities separately keeps the both-low case visible. The retired implementation is held out of the analysis path that produced the tables reported here.


\section{Experimental Setup}
\label{sec:setup}

Two question answering datasets are used, Natural Questions \citep{kwiatkowski-etal-2019} and HotpotQA \citep{yang-etal-2018}, at 1,000 queries each. Four embedding models index each corpus; retrieval takes the top five chunks of 256 tokens with 32-token overlap; three generators produce an answer from a byte-identical prompt at temperature 0 --- two proprietary (Claude Haiku 4.5 and GPT-4o-mini) and one open-weight run locally (Qwen2.5-7B-Instruct). This yields 4 x 2 x 3 x 1,000 = 24,000 generated answers, each re-scored by every evaluator. Prompt parity across generators is enforced by construction --- a single prompt builder, one user turn, no system message --- because a prompt difference would confound every cross-generator comparison in Section~\ref{sec:sensitivity}.

\begin{table}[t]
\centering
\resizebox{\columnwidth}{!}{%
\begin{tabular}{lrrr}
\toprule
\textbf{Embedding model} & \textbf{Paradigm} & \textbf{NDCG@5 (NQ)} & \textbf{NDCG@5 (HotpotQA)} \\
\midrule
all-mpnet-base-v2 & contrastive & 0.786 & 0.705 \\
BGE-M3 & multilingual & 0.794 & 0.809 \\
E5-large-instruct & instruction-tuned & 0.806 & 0.824 \\
text-embedding-3-small & proprietary & 0.830 & 0.767 \\
spread &  & 4.4 pts & 11.9 pts \\
\bottomrule
\end{tabular}
}
\caption{Retrieval quality of the four systems in the testbed. The ranking differs between datasets, and the spread is far larger on HotpotQA.}
\label{tab:1}
\end{table}

Three evaluators are compared. NLI-max is (1) with a DeBERTa-v3-large \citep{he-etal-2023} inference model, the common choice. Claim-min is (2) with the same inference model and sentence-level claim segmentation, chosen so that no second generator enters the measurement loop. AlignScore \citep{zha-etal-2023} is an independently trained alignment model. All three score identical inputs.

The answers are byte-identical across evaluators; the premises are not, and the difference is part of what the comparison measures. Premise construction follows each evaluator's intended interface. The entailment scorer receives each of the five retrieved chunks separately plus their concatenation truncated to the model's 512-token window, and the maximum is taken, since any one retrieved document may entail the answer and the concatenation exceeds the window. The claim-level scorer applies that same construction to each extracted claim. AlignScore receives the concatenated context as a single premise and performs its own splitting internally. Part of the difference reported in Section~\ref{sec:sensitivity} is therefore premise granularity rather than evaluator sensitivity; Section~\ref{sec:falsification} bounds that share by re-scoring the same cases with the same entailment model on the concatenated premise alone. The hypothesis is normalised identically for all three: markdown emphasis, headings and list markers are stripped before scoring. This matters because one generator emits markdown in 76.0\% of its attempted answers against 1.3\% and 0.3\% for the other two, and the entailment model responds to it --- see Section~\ref{sec:normalisation}.

One refusal rule is used throughout. An answer counts as a refusal when, after SQuAD-style normalisation, it begins with one of a fixed list of declining phrases; a single shared implementation applies it, so Sections V through VIII describe the same population. It excludes 5,776 of the 24,000 answers (24.1\%): 587 on NQ/Claude (14.7\%), 470 on NQ/Qwen (11.8\%), 1,162 on NQ/GPT-4o-mini (29.0\%), 1,166 on HotpotQA/Claude (29.1\%), 1,183 on HotpotQA/Qwen (29.6\%) and 1,208 on HotpotQA/GPT-4o-mini (30.2\%). All three generators prepend a "Based on the provided context," preamble to everything, refusals included, so the rule strips that preamble before the anchored test; without it 107 refusals, all of them Claude's, would be graded as attempts.


\section{How Sensitive Are These Evaluators to a Known Error?}
\label{sec:sensitivity}


\subsection{A controlled falsification}
\label{sec:falsification}

A faithfulness score is only useful if it responds when an answer stops being faithful. We construct that condition directly. For each answer containing a numeric value that appears in the retrieved context, we replace that value with one that does not appear in the context, preserving surface format, and re-score against the same chunks. Two eligibility rules are load-bearing: the original value must be present in the context, or the answer was never grounded in it; and the replacement must be absent, or the altered answer is accidentally supported and a low response is correct behaviour.

Cases are taken in dataset order rather than drawn at random: for each system, dataset and generator we use the first 200 eligible answers. The scored set is therefore a prefix of an already filtered population, and 4,784 cases are scored against 7,609 eligible across the grid --- 800 of 1,946 on NQ/Claude, 800 of 1,340 on NQ/Qwen, 800 of 882 on NQ/GPT-4o-mini, 800 of 1,551 on HotpotQA/Claude, 800 of 997 on HotpotQA/Qwen and 784 of 893 on HotpotQA/GPT-4o-mini. The two GPT-4o-mini cells are close to a census; the two Claude cells are sampled at two fifths and one half.

Two controls accompany it. An entity substitution replaces a grounded named entity instead of a value, giving a comparable edit. Scoring the untouched answer against another query's retrieved context gives the floor of each evaluator's scale on this data. 200 paired cases per cell, rebuilt deterministically so all three evaluators see identical inputs.

\begin{table*}[t]
\centering
\small
\begin{tabular}{lrrrrrr}
\toprule
\textbf{Dataset / generator} & \textbf{Evaluator} & \textbf{Orig} & \textbf{Falsified} & \textbf{Floor} & \textbf{Detected @0.5} & \textbf{Detected @floor gate} \\
\midrule
NQ / Claude Haiku 4.5 & NLI-max & 0.906 & 0.852 & 0.474 & 5.1\% [3.8-7.0] & 50.2\% [45.7-54.7] \\
 & Claim-min & 0.679 & 0.514 & 0.245 & 25.1\% [21.7-28.9] & 58.3\% [53.3-63.1] \\
 & AlignScore & 0.866 & 0.619 & 0.331 & 24.1\% [21.3-27.3] & 77.9\% [74.4-81.1] \\
NQ / Qwen2.5-7B & NLI-max & 0.836 & 0.574 & 0.470 & 35.1\% [31.7-38.8] & 71.0\% [66.1-75.4] \\
 & Claim-min & 0.693 & 0.397 & 0.378 & 48.1\% [44.0-52.2] & 79.0\% [74.6-82.8] \\
 & AlignScore & 0.863 & 0.428 & 0.240 & 51.8\% [48.2-55.4] & 75.5\% [72.1-78.7] \\
NQ / GPT-4o-mini & NLI-max & 0.887 & 0.436 & 0.506 & 56.6\% [52.9-60.1] & 85.2\% [81.9-88.1] \\
 & Claim-min & 0.863 & 0.390 & 0.490 & 61.0\% [57.4-64.6] & 91.0\% [88.1-93.2] \\
 & AlignScore & 0.912 & 0.272 & 0.160 & 75.4\% [72.2-78.4] & 82.1\% [79.2-84.7] \\
\midrule
HotpotQA / Claude & NLI-max & 0.702 & 0.618 & 0.427 & 13.9\% [11.3-17.0] & 49.7\% [42.3-57.2] \\
 & Claim-min & 0.370 & 0.255 & 0.161 & 34.3\% [29.0-40.0] & 48.9\% [41.8-56.1] \\
 & AlignScore & 0.786 & 0.454 & 0.475 & 48.9\% [45.2-52.5] & 93.7\% [90.7-95.8] \\
HotpotQA / Qwen2.5-7B & NLI-max & 0.547 & 0.277 & 0.245 & 54.8\% [50.1-59.5] & 69.6\% [64.0-74.7] \\
 & Claim-min & 0.424 & 0.160 & 0.178 & 66.8\% [61.6-71.6] & 76.9\% [71.3-81.7] \\
 & AlignScore & 0.749 & 0.255 & 0.304 & 74.5\% [71.1-77.7] & 96.6\% [94.4-98.0] \\
HotpotQA / GPT-4o-mini & NLI-max & 0.577 & 0.186 & 0.201 & 70.5\% [66.2-74.6] & 83.8\% [79.7-87.2] \\
 & Claim-min & 0.571 & 0.182 & 0.204 & 70.6\% [66.2-74.6] & 84.2\% [80.1-87.6] \\
 & AlignScore & 0.780 & 0.119 & 0.234 & 89.4\% [86.8-91.6] & 97.5\% [95.8-98.5] \\
\bottomrule
\end{tabular}
\caption{Response of three evaluators to a falsified value, on identical inputs, pooled over the four embedding systems (n = 800 per cell, 784 on HotpotQA/GPT-4o-mini). "Detected" is the share of answers that passed a threshold before falsification and fail it after, with Wilson 95\% intervals. "Floor" is the mean score of an untouched answer against an unrelated query's context; the floor gate is the 95th percentile of that same distribution, per cell and evaluator.}
\label{tab:2}
\end{table*}

On Claude's Natural Questions answers, NLI-max leaves the falsified answer at 0.852 against a floor of 0.474: the altered answer remains comfortably above any usable threshold. On GPT-4o-mini's answers the same edit lands at 0.436, essentially at that generator's floor of 0.506. The open-weight generator falls between them, at 0.574 against a floor of 0.470. The evaluator, the context and the perturbation code are identical; only the generator that wrote the answer differs.

That the third generator lands between the other two was predicted before it was run, from its assertion count alone (Section~\ref{sec:assertions}), and it holds in both datasets and under all three evaluators. The evaluator ordering is also stable: AlignScore is the most responsive, claim-min next and NLI-max least, in every one of the six cells. What differs across generators is not which evaluator responds most but how far apart the evaluators are.

Detection rates are only comparable across evaluators after the scale is normalised, because the evaluators do not share a floor, and we therefore report two gates. On this cell an untouched answer scored against an unrelated context averages 0.474 under NLI-max against 0.245 under claim-min, so a fixed 0.5 gate sits just above the floor for one evaluator and well above it for the other. The random-context distribution is also heavy-tailed --- its 95th percentile is 0.985 under NLI-max --- so an unrelated document clears 0.5 nearly half the time. The second gate is anchored at that 95th percentile, at which "detected" means the falsified answer is no longer better supported than an unrelated document, which is the same statement on every scale. Read that way NLI-max detects 50.2\% of falsifications on Claude's NQ answers rather than 5.1\%, and its gap to GPT-4o-mini narrows from roughly elevenfold to under twofold. The ordering of the three evaluators is unchanged at either gate. The floor-anchored gate is a comparison device and not an operating point: it sits above the mean untouched answer, so a deployed system could not use it.


\subsection{The effect is not verbatim copying}
\label{sec:copying}

A natural explanation is that one generator quotes the retrieved text back, so the evaluator scores a copy of its own premise and a single substituted value cannot move it. The data does not support this. Over the 4,784 falsification cases, the fraction of answer word 5-grams occurring in the retrieved context as a contiguous token sequence is 0.226 for Claude, 0.194 for Qwen and 0.235 for GPT-4o-mini. The explanation requires overlap to run opposite to sensitivity; it does not. The least sensitive generator does not copy the most, and the generator that copies least is the intermediate one rather than the most sensitive. The correlation between overlap and the response to falsification is +0.046 for Claude (95\% CI -0.004 to +0.096), +0.178 for Qwen (+0.127 to +0.229) and +0.166 for GPT-4o-mini (+0.118 to +0.212) --- indistinguishable from zero in the first case and positive in the other two, where the explanation requires all three to be negative.

The comparison also survives at matched overlap, which the correlation alone does not show. Grouping the cases into four overlap bands leaves Claude between +0.052 and +0.085, Qwen between +0.202 and +0.364 and GPT-4o-mini between +0.326 and +0.507. The ordering is the same in all four bands, including the band above 0.3 where all three generators are reusing a substantial part of the retrieved wording, and Claude is separated from GPT-4o-mini by at least fourfold throughout. Overlap is measured on tokens with punctuation mapped to spaces, so it is unaffected by the markdown normalisation described in Section~\ref{sec:normalisation}.


\subsection{Answers with more assertions are scored less sensitively}
\label{sec:assertions}

Segmenting each answer into sentence-level assertions and grouping by count gives a declining relationship, present in all three generators. It is not strictly monotone: the tail buckets are thin and each generator reverses direction once within them.

\begin{table*}[t]
\centering
\small
\resizebox{\textwidth}{!}{%
\begin{tabular}{lrrrrrrrrr}
\toprule
 & \multicolumn{3}{c}{\textbf{Claude Haiku 4.5}} & \multicolumn{3}{c}{\textbf{Qwen2.5-7B}} & \multicolumn{3}{c}{\textbf{GPT-4o-mini}} \\
\cmidrule(lr){2-4}\cmidrule(lr){5-7}\cmidrule(lr){8-10}
\textbf{Assertions} & \textbf{n} & \textbf{NLI-max} & \textbf{Claim-min} & \textbf{n} & \textbf{NLI-max} & \textbf{Claim-min} & \textbf{n} & \textbf{NLI-max} & \textbf{Claim-min} \\
\midrule
1 & 91 & +0.102 & +0.102 & 413 & +0.409 & +0.409 & 680 & +0.485 & +0.485 \\
2 & 300 & +0.076 & +0.151 & 205 & +0.152 & +0.219 & 77 & +0.360 & +0.467 \\
3 & 194 & +0.032 & +0.229 & 103 & +0.051 & +0.150 & 26 & +0.041 & +0.401 \\
4 & 103 & +0.024 & +0.162 & 40 & +0.024 & +0.076 & 10 & +0.204 & +0.192 \\
5 or more & 112 & +0.024 & +0.146 & 39 & +0.072 & +0.126 & 7 & +0.032 & +0.002 \\
\bottomrule
\end{tabular}
}
\caption{Drop in score when one grounded value is falsified, by number of assertions in the answer, on Natural Questions. Under NLI-max the response decays as answers accumulate assertions; under claim-level aggregation it does not. The GPT-4o-mini rows beyond two assertions rest on 26, 10 and 7 cases and should be read as such. In the single-assertion row the two aggregates coincide exactly, which is the degenerate case they must reduce to.}
\label{tab:3}
\end{table*}

The three generators average 3.08, 1.77 and 1.25 assertions per attempted answer on this dataset, with median answer lengths of 55, 25 and 15 words. This largely accounts for the cross-generator difference in Table II: a falsified value is a progressively smaller edit to the hypothesis as correct surrounding text accumulates, and the entailment judgement is dominated by the remainder. It does not account for it entirely --- within the single-assertion bucket the three generators still differ, +0.102 against +0.409 and +0.485. As a check on the claim-level implementation, every single-assertion answer scores identically under both aggregates (n = 2,472, largest absolute difference 1.7e-05, which is batch-padding noise), the expected degenerate case. That check is what surfaced the normalisation defect reported in Section~\ref{sec:normalisation}: before markdown was stripped for both aggregates it failed on every one of the 131 formatted answers, by as much as 0.729.

Claim-level aggregation is a partial remedy rather than a repair. Expressed as a fraction of the distance to the random-context floor, the median response on Claude's answers moves from 0.01 to 0.03, and the minimum moves in the correct direction in 55\% of cases: when the falsified assertion is not already the weakest, the minimum does not move at all. A second dilution remains within assertions --- at matched assertion count the three generators still differ (+0.102, +0.409 and +0.485), and their single-assertion answers average 23.9, 19.5 and 16.1 words.

The third generator is what makes this residual legible. In aggregate it sits between the other two, but conditioned on assertion count it does not: at one assertion it is +0.409 against GPT-4o-mini's +0.485 and Claude's +0.102. Its intermediate aggregate is produced by its assertion \emph{distribution} rather than by an intermediate per-assertion sensitivity. The remaining gap is therefore specific to one generator and not a smooth function of verbosity, which two generators cannot distinguish and three can.


\section{The Evaluation Population Can Manufacture a Relationship}
\label{sec:population}

All three generators decline to answer on a substantial share of queries, and the rate depends on retrieval quality: on HotpotQA the weakest retriever draws 0.378 to 0.398 refusals and the strongest 0.224 to 0.241, in every generator. A refusal is not an attempt at a grounded answer, so including refusals transfers a retrieval effect into the faithfulness average, at no cost in plausibility.

\begin{table}[t]
\centering
\resizebox{\columnwidth}{!}{%
\begin{tabular}{llrrrrr}
\toprule
\textbf{HotpotQA} & & \textbf{all-mpnet} & \textbf{text-emb-3-small} & \textbf{BGE-M3} & \textbf{E5-instruct} & \textbf{spread} \\
\midrule
& NDCG@5 & 0.705 & 0.767 & 0.809 & 0.824 & \\
\midrule
Claude & refusal rate & 0.393 & 0.301 & 0.248 & 0.224 & \\
& faithfulness, pooled & 0.538 & 0.589 & 0.609 & 0.622 & 0.085 \\
& faithfulness, answered & 0.710 & 0.709 & 0.693 & 0.712 & 0.018 \\
\midrule
Qwen & refusal rate & 0.378 & 0.302 & 0.262 & 0.241 & \\
& faithfulness, pooled & 0.694 & 0.664 & 0.641 & 0.647 & 0.053 \\
& faithfulness, answered & 0.568 & 0.569 & 0.569 & 0.566 & 0.003 \\
\midrule
GPT-4o-mini & refusal rate & 0.398 & 0.323 & 0.255 & 0.232 & \\
& faithfulness, pooled & 0.766 & 0.731 & 0.694 & 0.716 & 0.072 \\
& faithfulness, answered & 0.651 & 0.636 & 0.622 & 0.649 & 0.029 \\
\bottomrule
\end{tabular}
}
\caption{Pooling refusals into the faithfulness average produces a column that tracks retrieval quality in all three generators, with a spread of 0.053 to 0.085 --- but in opposite directions: positively for Claude (Spearman +0.99 against NDCG@5) and negatively for the other two (-0.98 and -0.92). Restricted to answered rows the spread collapses to 0.003-0.029 and the ordering follows retrieval quality in none of them. Systems are ordered by NDCG@5.}
\label{tab:4}
\end{table}

Any evaluation population whose membership depends on the system under comparison can transfer an effect from the selection into the measurement. The weakest retriever refuses most often, so its pooled average is the one most displaced by whatever score refusals receive.

What that score is turns out to be a property of the generator's refusal wording rather than of refusing. Under NLI-max, Claude's HotpotQA refusals average 0.272 to 0.353 against 0.693 to 0.712 for its attempted answers, so pooling drags the weak retriever down and the pooled column ascends with NDCG@5. Qwen's refusals average 0.846 to 0.903 and GPT-4o-mini's 0.905 to 0.940, in both cases \emph{above} their attempted answers, so pooling lifts the weak retriever and the same column descends. Both readings are available from the same experiment and neither is about faithfulness: the pooled column is close to a monotone function of the refusal rate, and its sign is set by how the entailment model happens to score one generator's standard refusal sentence. This is the sharper version of the point --- a defender of the pooled statistic cannot appeal to a consistent direction of bias, because the direction reverses between generators on identical queries and identical retrieval.

All faithfulness results in this paper exclude refusals and are paired on the queries every system answered.


\subsection{Membership Is Set by a Grader, and Graders Disagree}
\label{sec:graders}

The remainder of this section rests on a judge model, which was run on the two proprietary arms only; it and Tables V and VI therefore cover sixteen of the twenty-four cells, and no claim here is made about the open-weight generator.

Excluding refusals requires deciding which answers are refusals, and that decision is made by a rule rather than given by the data. We compare two defensible rules: a string-matching heuristic keyed to the refusal templates the prompts offer, and a judge model shown the question, the reference answers and the answer. They agree on the large majority of rows and disagree consistently at the margin.

\begin{table}[t]
\centering
\resizebox{\columnwidth}{!}{%
\begin{tabular}{lrrrr}
\toprule
\textbf{HotpotQA, Claude} & \textbf{all-mpnet} & \textbf{text-emb-3-small} & \textbf{BGE-M3} & \textbf{E5-instruct} \\
\midrule
Refusal rate, heuristic & 0.376 & 0.279 & 0.224 & 0.207 \\
Refusal rate, judge & 0.422 & 0.332 & 0.277 & 0.252 \\
\bottomrule
\end{tabular}
}
\caption{Two refusal detectors on the same answers. The judge finds roughly four to five more refusals per hundred answers in every column, so the disagreement is a level shift rather than noise, and it does not reorder the systems.}
\label{tab:5}
\end{table}

A uniform shift that preserves the ordering would ordinarily be unremarkable. It is not unremarkable here, because the excluded rows are not a random sample of the population: they are the rows nearest the boundary between refusing and attempting. Substituting one detector for the other moves nineteen rows on Natural Questions with Claude and fifty-five on HotpotQA with Claude, and changes the significance verdict in two of the eight dataset-generator-evaluator cells it was run over. Under the entailment aggregate, Natural Questions with Claude moves from no detectable difference (p = 0.223) to a detectable one (p = 0.048); under AlignScore, HotpotQA with Claude moves from p = 0.151 to p = 0.049. Both land within a thousandth of the conventional threshold, which is the point: the verdict in these cells is not a property of the systems.

The count is more stable than the cells: the number of dataset-generator cells whose verdict changes between evaluators is unchanged by substituting one detector for the other, but the cells are not the same ones. We therefore report the count as the finding and treat any individual cell as undetermined.


\subsection{The Correctness Grader Admits Refusals to the Wrong Population}
\label{sec:refusals}

The same problem appears once faithfulness is conditioned on correctness. String containment grades an answer correct when a reference answer string occurs in it, and a short reference string can occur inside a refusal. On a 200-answer calibration, ten refusals were graded correct by containment and the judge overturned nine of them. A refusal admitted to the correct population carries a near-floor faithfulness score into the correct-answer average, which lowers the baseline that any correct-versus-incorrect comparison is measured against. Under graded judgement no refusal is graded correct in any of the sixteen judged cells.

The effect is large enough to change a sign. Measured with containment and a baseline that excluded refusals from the incorrect side only, answers on Natural Questions with Claude appeared more grounded when wrong than when right, by 0.018 to 0.050 across the four systems. Excluding refusals from both sides reduces this to between -0.006 and +0.018. Grading correctness with the judge instead of containment moves it to between -0.051 and -0.113, in the opposite direction. Each of the three steps is a correction to the grader rather than a change to the answers, and all three move the quantity the same way.

Tested directly, no positive value of this quantity survives. Across all sixteen judged cells a two-sample permutation test with Holm correction finds three significant, all negative, and none of the eight positive-signed cells has a confidence interval excluding zero. The clearest single result is E5-instruct on Natural Questions, where the same effect appears under both judged generators (-0.113 and -0.115). We report no evidence that answers are more grounded when they are wrong. The intervals are wide, between 0.08 and 0.13, because only 44 to 84 wrong answered rows remain per cell once refusals are excluded; these tests bear on the general claim and not on individual systems.


\subsection{A Contingency Cell That Mostly Counts Refusals}
\label{sec:cell}

The same accounting governs the contingency between retrieval success and correctness. The cell in which retrieval succeeded and the answer was still wrong is the one that bears on whether retrieval quality is sufficient, and most of it consists of answers that were never attempted.

\begin{table}[t]
\centering
\resizebox{\columnwidth}{!}{%
\begin{tabular}{lrrrr}
\toprule
\textbf{Retrieval hit and answer incorrect} & \textbf{Claude, NQ} & \textbf{Claude, HQA} & \textbf{GPT, NQ} & \textbf{GPT, HQA} \\
\midrule
Share of all queries & 17.4\% & 34.0\% & 31.7\% & 36.3\% \\
Of that cell, refusals & 69.3\% & 86.8\% & 79.0\% & 78.1\% \\
Attempted and wrong & 5.4\% & 4.5\% & 6.7\% & 7.9\% \\
\bottomrule
\end{tabular}
}
\caption{Correctness from graded judgement. The raw cell is between 17 and 36 per cent of queries; the share that is a grounded, committed, wrong answer is between 4.5 and 7.9 per cent.}
\label{tab:6}
\end{table}

Read without the decomposition, the raw cell would support the claim that good retrieval frequently fails to produce a correct answer. Most of what it counts is a system declining to answer, and refusal rates differ by more than twenty points across the systems compared. The residual is a real phenomenon and it is roughly a fifth the size the raw cell implies.


\section{The Verdict Depends on the Evaluator}
\label{sec:verdict}

We now compare the four retrieval systems on faithfulness, refusals excluded, paired on a common query subset, with a paired bootstrap of 10,000 resamples.

\begin{table*}[t]
\centering
\small
\begin{tabular}{llrrrr}
\toprule
\textbf{Dataset / generator} & \textbf{Evaluator} & \textbf{n} & \textbf{Spread} & \textbf{p} & \textbf{p (Holm)} \\
\midrule
NQ / Claude & NLI-max & 784 & 0.0107 & 0.155 & 1.000 \\
 & Claim-min & 784 & 0.0216 & 0.104 & 0.978 \\
 & AlignScore & 784 & 0.0114 & 0.0081 & 0.122 \\
NQ / Qwen & NLI-max & 790 & 0.0055 & 0.538 & 1.000 \\
 & Claim-min & 790 & 0.0076 & 0.509 & 1.000 \\
 & AlignScore & 790 & 0.0072 & 0.204 & 1.000 \\
NQ / GPT-4o-mini & NLI-max & 628 & 0.0130 & 0.098 & 0.978 \\
 & Claim-min & 628 & 0.0151 & 0.075 & 0.823 \\
 & AlignScore & 628 & 0.0080 & 0.117 & 0.978 \\
\midrule
HotpotQA / Claude & NLI-max & 517 & 0.0232 & 0.159 & 1.000 \\
 & Claim-min & 517 & 0.0345 & 0.053 & 0.739 \\
 & AlignScore & 517 & 0.0164 & 0.063 & 0.823 \\
HotpotQA / Qwen & NLI-max & 499 & 0.0123 & 0.568 & 1.000 \\
 & Claim-min & 499 & 0.0181 & 0.294 & 1.000 \\
 & \textbf{AlignScore} & 499 & \textbf{0.0406} & \textbf{0.0011} & \textbf{0.0187} \\
HotpotQA / GPT-4o-mini & NLI-max & 510 & 0.0267 & 0.065 & 0.823 \\
 & Claim-min & 510 & 0.0378 & 0.0070 & 0.112 \\
 & \textbf{AlignScore} & 510 & \textbf{0.0461} & \textbf{0.0005} & \textbf{0.0090} \\
\bottomrule
\end{tabular}
\caption{The same comparison under three evaluators, corrected over the eighteen comparisons in the table. Two of six cells change verdict at alpha = 0.05 and survive correction, both under AlignScore and both on HotpotQA; two further comparisons are significant uncorrected and not after correction. Queries, answers, refusal handling and test are identical throughout. With 10,000 resamples the smallest attainable non-zero p is 1e-4, so the cell at p = 0.0005 is 5 extreme resamples of 10,000 and not a clamped floor.}
\label{tab:7}
\end{table*}

The statistic is the spread: the difference in mean faithfulness between the best and the worst of the four systems, with the p-value from a paired bootstrap over that pair, 10,000 resamples, re-centred at zero \citep{dror-etal-2018}. The range is the quantity the decision this table informs is exposed to --- a practitioner adopts one embedding model, and what they stand to lose is the distance to the best alternative, not the dispersion of the set. It is however a summary of two systems selected for being the extremes, so the p-value beside it is not adjusted for that within-cell selection, and the two intermediate systems do not appear in it at all. We therefore treat the range as descriptive, and report the direction across all four systems separately in Table VIII rather than reading it off the extremes. An omnibus test over the four systems would be the stricter statistic; we do not run one here.

Under the common choice of evaluator, no cell shows a detectable difference in any of the six. Under the evaluator measured in Section~\ref{sec:sensitivity} to be the most responsive to known errors, two do. A result reported only under the first evaluator would be an artifact of instrument resolution rather than a property of the systems. The two surviving cells are both on HotpotQA, the dataset with the wider retrieval spread, and neither is on Claude, the generator whose answers the entailment aggregate is least able to resolve.

Adding a third generator made the correction stricter --- eighteen comparisons rather than twelve --- and the count rose anyway, which the design does not guarantee. It also produced the strongest single result here: under AlignScore on HotpotQA, all three generators rank the four retrieval systems identically (all-mpnet $<$ text-embedding-3-small $<$ BGE-M3 $<$ E5-instruct), while under NLI-max no two of them agree. This should be read carefully. The generators are independent of one another, but they answer from the \emph{same} retrieved contexts, so this is three generators agreeing about a property of the shared retrieval rather than three independent replications of a result. What it does establish is that one evaluator resolves that property consistently across generators and another does not, on identical inputs.

That count of two is itself not invariant to the correction. Two comparisons sit either side of alpha = 0.05 depending on which comparisons are treated as one family. NQ/Claude under AlignScore is at p = 0.0081 uncorrected and is adjusted to 0.122 over the eighteen in the table, 0.032 over the six sharing an evaluator and 0.024 over the three in its own cell; HotpotQA/GPT-4o-mini under claim-min is at p = 0.0070 and is adjusted to 0.112, 0.042 and 0.014 respectively. Either smaller family would make the count four rather than two. Correcting within evaluator is defensible, since the six NLI tests establish which cells are null and so precede the question rather than competing with it. We report the eighteen-comparison family because it is the most conservative of the three and because it was fixed before the correction was applied; adopting a smaller family after observing that it restores cells would be a selection we could not account for. We state the alternative rather than only the choice, because a count that moves with the multiplicity family is the same kind of dependence on an analysis decision that this paper reports for the choice of evaluator.


\subsection{Direction is dataset-dependent}
\label{sec:direction}

Where a difference exists, it does not point the same way. Spearman correlation between NDCG@5 and faithfulness across the four systems, computed in every generator-evaluator combination:

\begin{table}[t]
\centering
\resizebox{\columnwidth}{!}{%
\begin{tabular}{llrr}
\toprule
\textbf{Generator} & \textbf{Evaluator} & \textbf{Natural Questions} & \textbf{HotpotQA} \\
\midrule
Claude & NLI-max & -1.00 & +0.20 \\
 & Claim-min & -0.60 & +0.60 \\
 & AlignScore & -0.20 & +1.00 \\
Qwen & NLI-max & -0.80 & +0.60 \\
 & Claim-min & +0.63 & -0.40 \\
 & AlignScore & -0.40 & +1.00 \\
GPT-4o-mini & NLI-max & -0.80 & +0.80 \\
 & Claim-min & -0.80 & +0.80 \\
 & AlignScore & 0.00 & +1.00 \\
\bottomrule
\end{tabular}
}
\caption{With four systems these coefficients cannot carry significance individually. Under NLI-max and AlignScore the pattern is clean in all twelve combinations: every HotpotQA coefficient is positive and every Natural Questions coefficient is zero or negative. Under claim-min it fails in two of six, both on the open-weight generator. The one non-canonical value, +0.63, arises from a tie between two systems.}
\label{tab:8}
\end{table}

On the multi-hop dataset, better retrieval accompanies more faithful answers; on the single-hop dataset it does not. A plausible reading is that composing an answer from several retrieved pieces makes grounding sensitive to how much of the needed evidence is present, while a single-fact answer is not. We do not test this, and note that the experiment varies embedding model rather than manipulating retrieval quality directly, so the correlation should not be read as an effect of NDCG@5 itself.

The qualification is that the pattern is itself evaluator-dependent, which is the paper's subject applied to its own subsidiary result. It is exact under NLI-max and AlignScore and breaks under claim-min on the open-weight generator, in both datasets and in opposite directions. A two-evaluator or two-generator version of this table would have shown a clean dataset effect; the third of each is what reveals that the direction is not a property of the dataset alone.


\subsection{A normalisation defect that produced a result}
\label{sec:normalisation}

An earlier version of this analysis, run before the open-weight generator was added, reported two of the four cells then in the design rather than one. The difference was not a change of test, data or population but of text normalisation, and it is worth reporting because it is an instance of the paper's own subject.

The entailment scorer originally read the answer exactly as the generator wrote it, while the claim-level scorer stripped markdown from each claim before scoring. The two aggregates were therefore not reading the same hypothesis. The discrepancy is invisible in aggregate and was found by the degenerate-case check in Section~\ref{sec:assertions}: with a single claim and no markdown the two must be numerically identical, and they were, to 1.7e-05; with markdown they differed on every one of the 131 cases then in the grid, by up to 0.729 on a pair of asterisks alone. Markdown is not evenly distributed --- it appears in 76.0\% of Claude's attempted answers against 1.3\% and 0.3\% for the other two generators --- so the formatting habits of one generator were entering a comparison between embedding systems.

Stripping markdown for all three evaluators removes the asymmetry and restores the identity on every single-assertion answer. It also removed the significance of the NQ/Claude cell under AlignScore: its spread fell from 0.0179 to 0.0114, and it is the cell that Table VII still places just outside the threshold. That cell's apparent evaluator-dependence was substantially an artifact of formatting, and the count reported above is the corrected one. Neither cell that survives correction in Table VII is affected: both are on generators that emit almost no markdown. We record it rather than silently reporting the smaller number because the failure mode generalises: a preprocessing choice applied to one aggregate and not another is invisible in every summary statistic, survives every significance test, and is detectable only by an identity that the two aggregates must satisfy. Metrics that admit such an identity should be checked against it.


\section{What an Equivalence Claim Requires}
\label{sec:equivalence}

Four of the six cells in Table VII show no detectable difference under any of the three evaluators. Absence of a detected difference is not evidence of equivalence, so we test equivalence directly with two one-sided tests \citep{schuirmann-1987}. A TOST result depends entirely on the margin, and a margin is a substantive judgement rather than a statistical fact --- so we report a curve rather than assert a point.

\begin{table*}[t]
\centering
\small
\begin{tabular}{lrrrrrr}
\toprule
\textbf{Dataset / generator} & \textbf{Evaluator} & \textbf{$\pm$0.01} & \textbf{$\pm$0.02} & \textbf{$\pm$0.03} & \textbf{$\pm$0.05} & \textbf{$\pm$0.10} \\
\midrule
NQ / Claude & NLI-max & 0/6 & 4/6 & 6/6 & 6/6 & 6/6 \\
 & Claim-min & 0/6 & 0/6 & 2/6 & 6/6 & 6/6 \\
 & AlignScore & 3/6 & 6/6 & 6/6 & 6/6 & 6/6 \\
NQ / Qwen & NLI-max & 0/6 & 3/6 & 6/6 & 6/6 & 6/6 \\
 & Claim-min & 0/6 & 0/6 & 6/6 & 6/6 & 6/6 \\
 & AlignScore & 1/6 & 6/6 & 6/6 & 6/6 & 6/6 \\
NQ / GPT-4o-mini & NLI-max & 0/6 & 3/6 & 6/6 & 6/6 & 6/6 \\
 & Claim-min & 0/6 & 3/6 & 6/6 & 6/6 & 6/6 \\
 & AlignScore & 1/6 & 6/6 & 6/6 & 6/6 & 6/6 \\
\midrule
HotpotQA / Claude & NLI-max & 0/6 & 0/6 & 0/6 & 5/6 & 6/6 \\
 & Claim-min & 0/6 & 0/6 & 0/6 & 3/6 & 6/6 \\
 & AlignScore & 0/6 & 2/6 & 5/6 & 6/6 & 6/6 \\
HotpotQA / Qwen & NLI-max & 0/6 & 0/6 & 0/6 & 6/6 & 6/6 \\
 & Claim-min & 0/6 & 0/6 & 0/6 & 6/6 & 6/6 \\
 & AlignScore & 0/6 & 0/6 & 0/6 & 4/6 & 6/6 \\
HotpotQA / GPT-4o-mini & NLI-max & 0/6 & 0/6 & 1/6 & 4/6 & 6/6 \\
 & Claim-min & 0/6 & 0/6 & 1/6 & 4/6 & 6/6 \\
 & AlignScore & 0/6 & 0/6 & 0/6 & 4/6 & 6/6 \\
\bottomrule
\end{tabular}
\caption{System pairs judged equivalent, of six tested, across margins. A claim of equivalence is a claim about a particular column --- and about a particular row: at $\pm$0.03 on NQ/Claude, NLI-max calls all six pairs equivalent and claim-min two, on the same answers.}
\label{tab:9}
\end{table*}

Section~\ref{sec:sensitivity} supplies an external anchor for choosing that column. Falsifying one grounded value moves NLI-max by 0.049 to 0.095 on Claude's answers. A margin of $\pm$0.05 on that evaluator is therefore approximately the size of one fabricated fact, which is not a negligible difference. Stating equivalence at $\pm$0.05 without that context asserts more than the instrument supports.

The anchor is evaluator-specific, and so is the verdict it licenses. The same margin corresponds to a different number of fabricated facts on each evaluator --- one falsified value moves claim-min by 0.165 and AlignScore by 0.247 on the same Claude answers --- and the equivalence counts move with it: at $\pm$0.03 on NQ/Claude, NLI-max calls all six pairs equivalent and claim-min two, from the same generations. An equivalence claim therefore has to name its evaluator for the same reason a difference claim does, and additionally has to say what its margin is worth on that evaluator's scale.


\section{Discussion}
\label{sec:discussion}

For work that reports RAG faithfulness, we suggest the following.

\begin{itemize}
\item Name the evaluator in the claim. "Faithfulness did not differ" is incomplete; two of our six cells reverse between evaluators on identical data.
\item Report the evaluation population explicitly, and whether membership depends on the system being compared. Pooling refusals produced a column tracking retrieval quality in all three generators here, in opposite directions depending on how the evaluator scores that generator's refusal wording (Table IV).
\item Prefer claim-level aggregation to a maximum over the whole answer, and expect it to help most where answers carry several assertions.
\item Calibrate the margin against a measured effect before claiming equivalence. A falsification experiment of the kind in Section~\ref{sec:sensitivity} costs nothing beyond re-scoring.
\end{itemize}

For a document question answering system, the practical reading of Table II is that a grounding threshold on a verbose generator detects a minority of falsified values at a fixed gate: 5.1\% under the common evaluator, 25.1\% under the most responsive one tested. A deterministic check that every value in an answer appears in the retrieved context addresses this class of error directly and is complementary to an entailment score. Run over the same cases as Table II, it recalls every falsified value in all six cells at a false-positive rate of 0.8\% to 5.2\% on untouched answers, against 5.1\% detection for NLI-max on the hardest of those cells. Two qualifications matter. Numbers that are not assertions about the world --- the chunk indices the prompt supplies and the markers of an ordered list --- are neither asserted of the world nor expected in the retrieved text, and counting them raises the false-positive rate on Claude's HotpotQA answers from 5.2\% to about 39\%. And the residual false positives are quantities the answer derives rather than copies: counts, sums, elapsed years, formats converted. A literal check flags all of those, so it is a complement to an entailment score and not a replacement for one.

Checking symbolic content against the source is not a new proposal, and we do not claim it as one. \citet{goodrich-etal-2019} score generated text by extracting subject-relation-object triples and comparing them against the source document. \citet{nan-etal-2021} define entity-level consistency metrics for summarisation and evaluate them beside entailment-based alternatives. \citet{zhao-etal-2020} address quantity hallucination directly, verifying dates, numbers and monetary amounts in generated summaries against the source. That numbers are a weak point is also documented at the level of the representations these systems are built on: \citet{wallace-etal-2019} find number magnitude encoded unevenly across standard embeddings and least reliably by subword models. Entailment is separately a loose proxy for attribution --- the AIS framework \citep{rashkin-etal-2023} distinguishes whether a statement is supported by an identified source from whether it is plausible or entailed in a general sense, which is the distinction a structurally similar but differently attributed claim exploits.

What this paper adds is narrower than the check itself: the controlled minimal pair. A single numeric value is replaced in an otherwise byte-identical generation, and the same case is scored by all three evaluators, so the quantity measured is each evaluator's response to a known error on identical input rather than its agreement with a human label on some other corpus. The scope of that measurement is numeric values, three generators --- two proprietary and one open-weight --- and two datasets.


\section{Conclusion}
\label{sec:conclusion}

Holding a 24,000-answer generation grid fixed and varying only the faithfulness evaluator changes the conclusion about whether embedding choice affects faithfulness in two of six dataset-generator cells. The differences between evaluators are measurable directly: a falsified grounded value moves the common entailment aggregate by 0.054 on one generator, 0.262 on a second and 0.451 on a third, and that gap is largely accounted for by how many assertions an answer contains rather than by the retrieval system. Reported without naming the evaluator, the evaluation population, the aggregation and the equivalence margin, a faithfulness result is not reproducible in the sense that matters --- another group applying a different standard choice to the same generations would reach a different conclusion.

\section*{Limitations}

The claim is bounded by the grid it was measured on and by the graders that define its populations.

\begin{itemize}
\item Four embedding models, two English datasets and three generators, two proprietary and one open-weight. The three span 1.12 to 2.73 assertions per answer, which is the axis the mechanism in Section~\ref{sec:assertions} runs along, but they remain three points on it and two of the three are instruction-tuned chat models from the same broad family of training recipes.
\item Correctness is graded by a judge model rather than by string overlap. Exact match is 0.000 across all 24,000 answers because the generators do not emit bare answer spans. String containment, the obvious fallback, is not a bound in either direction: against the judge it understates accuracy by 5 to 15 points in twelve cells and overstates it by 3 to 7 points in the four HotpotQA cells under Claude, where long answers mention a reference string while asserting something else. The judge is itself a language model applied to its own outputs, and its agreement with a human annotator on this data has not been established.
\item Retrieval hit is defined as the gold answer string appearing in the retrieved context. This is a strong definition, and it does not separate retrieval failure from failures of composition, entity resolution or inference once the evidence is present.
\item The Natural Questions sample retains only queries whose annotated answer falls inside the retained context window. This makes the retrieval task more constrained than open-domain RAG, where evidence may be distributed, partial or implicit.
\item Pairing on the queries every system answered removes exactly those queries where retrieval quality determined whether an answer was produced. The reported comparison is therefore conditional on that subset, which is narrower than an unconditional comparison.
\item Effect sizes are small where detected --- 0.046 and 0.041 --- against retrieval spreads of 4.4 and 11.9 NDCG@5 points. These are detectable differences, not large ones.
\item The three evaluators receive byte-identical answers but not identical premises. The entailment scorer is given each of the five retrieved chunks separately, plus their concatenation truncated to the model's window, and the maximum is taken; the claim-level scorer applies that construction per extracted claim; AlignScore is given the concatenated context as one premise and performs its own splitting. Each is the evaluator used as its authors intend, but it means part of the difference reported in Section~\ref{sec:sensitivity} is premise granularity rather than evaluator sensitivity. Re-scoring the same cases with the same entailment model on the concatenated premise bounds that share between none and half of the gap depending on the cell, and at about a quarter of it on Claude's answers. The hypothesis, by contrast, is now normalised identically for all three (Section~\ref{sec:normalisation}); it was not in an earlier version, and that difference alone decided one cell.
\item Which answers count as refusals is fixed by a rule rather than given by the data, and that rule sets the evaluation population for every result in Sections V through VIII. It is therefore a researcher degree of freedom as well as a finding: substituting one defensible detector for another moves the rows nearest the boundary between refusing and attempting, and carries the uncorrected p-value across the conventional threshold in two of eight cells (Section~\ref{sec:graders}). Those two do not survive correction for multiple comparisons, but the population every reported number is computed on still depends on the choice. The rule used here strips the generators' standard preamble before matching; without that step 107 refusals, all from one generator, would be graded as attempts.
\item The falsification cases are the first 200 eligible answers in dataset order for each system, dataset and generator, rather than a random sample of the eligible answers: 4,784 cases scored against 7,609 eligible across the grid. Eligibility requires a grounded numeric value in the answer together with a replacement value absent from the context, so the scored set is a prefix of an already filtered population. Cells whose eligible count is near 200 are close to a census; the largest cells are sampled at about two fifths.
\item The repository configures seven embedding models and this paper reports four: all-mpnet-base-v2 (contrastive), BGE-M3 (multilingual), E5-large-instruct (instruction-tuned) and text-embedding-3-small (proprietary), chosen to span pretraining paradigms and to include one closed-source arm. GTE-large, Instructor-XL and jina-embeddings-v3 \citep{sturua-etal-2024} are configured and runnable but were not run at this scale: each additional model is a further generation pass over both datasets and all three generators, and the grid is a testbed for the evaluator comparison rather than a survey of embedding models.
\item Correctness is graded only on the two proprietary arms. Section~\ref{sec:graders} and Tables V and VI therefore cover sixteen of the twenty-four cells, and the open-weight generator enters no result that conditions on correctness. Nothing in Sections~\ref{sec:sensitivity}, \ref{sec:verdict} or \ref{sec:equivalence} depends on that grading.

\item The correct-versus-incorrect faithfulness comparison rests on 44 to 84 attempted wrong answers per cell. Excluding refusals and grading correctness with a judge both reduce that population, and the resulting confidence intervals are 0.08 to 0.13 wide. The tests reported in Section~\ref{sec:refusals} bear on the general claim and do not resolve individual systems.
\end{itemize}

\bibliography{custom}

\end{document}
